% page 161 du fac-simile (image p-160 du scan)
% bloc 172.7 x 233.3 mm ; marge gauche 29.3 mm
\pgc{7.620mm}{0.000mm}{
\l{}
\l{}
\l{\cel{58}151}
\l{}
\l{}
\l{}
\l{existar.\cel{2}(netrans.) Apartenar a la nuna realaji. - DEFIS.}
\l{}
\l{exkavar.\cel{2}(\sou{trans.}) Facar, kavigante. - EFIS.}
\l{}
\l{\sou{exk}luzar. (trans.)\cel{2}I. Pozar (ulu) exter kozo kom ne esante plu}
\l{\cel{3}partoprenonta olu. - II. Retropulsar (kozo) kom nekonciliebla}
\l{\cel{3}kun altro. - DEFIRS.}
\l{}
\l{\sou{exkomunikar}. (\sou{trans.)}\cel{2}(\sou{religio kristana}). Forjetar (ulu) de la}
\l{\cel{3}komuniono (katolika). - DEFIRS.}
\l{}
\l{\sou{exkoriar}. (\sou{trans.)}\cel{2}(\sou{medic.)}\cel{3}Skrachar surfacale, entamante nur}
\l{\cel{3}la pelo. - EFIS.}
\l{}
\l{\sou{exkremento}. (\sou{fiziol.}) Materio liquida o solida (sudoro, muko nazala,}
\l{\cel{3}urino, feko) quan ula organi ekpulsas de la korpo. DEFIRS.}
\l{}
\l{\sou{exkursar}.\cel{1}\sou{(netrans.}) Cirkular explorante ula quanto de la surfaco,}
\l{\cel{3}de la areo di lando. - DEFIRS.}
\l{}
\l{\sou{exkuzar}. (\sou{trans.)}\cel{2}Alegar,favore di ulu,motivi por mingravigar to}
\l{\cel{3}quon onu reprochas a lu, o por justifikar lu pri lo. -DEFIRS.}
\l{}
\l{\sou{exodo}.\cel{2}I. (\sou{pri la Hebrei antiqua}).. Libro duesma di Biblo, qua}
\l{\cel{3}kontenas la historio di la ekiro de Egiptia. - II. (\sou{metaf.})}
\l{\cel{3}Ekmigro grand-amase de populo. - DEFIS.}
\l{}
\l{\sou{exogena}.\cel{1}\sou{(bot.}) Di qua la kresko eventas adexter su ipsa, per}
\l{\cel{3}la strati extera.}
\l{}
\l{\sou{exorcisar}.\cel{2}(\sou{trans.)}\cel{2}(\sou{teol. kristana}). Agar a la demoni (por expul-}
\l{\cel{3}sar li de la korpo di demonozi), od a la demonozo, la pregi, la}
\l{\cel{3}ceremonii di la eklezio katolika. - DEFIS.}
\l{}
\l{\sou{exordiar}. (\sou{trans.})\cel{2}(\sou{retoriko}). Pri diskurso : komencar, preparante}
\l{\cel{3}la atenco e la afableso-sentimenti. - DEFIS.}
\l{}
\l{\sou{exostoso}.\cel{2}I. (\sou{medic.)}\cel{2}Surkreskajo morbala sur la surfaco od}
\l{\cel{3}interne di osto. - II. (\sou{bot.}) Surkreskajo sur la trunko, la}
\l{\cel{3}branchi di ula arbori. - DEF.}
\l{}
\l{\sou{exotera}. I. (\sou{filoz}.) Dicesis pri la dicipuli qui vartis sua inicieso}
\l{\cel{3}pri ula doktrini neprofunda - pri ula libri kompozita tale ke}
\l{\cel{3}lia kontenajo esis komprenebla da la plebeyo. - II. (\sou{okultismo}).}
\l{\cel{3}Dicesas pri la docajo quan onu darfas divulgar ad omna homi.-}
\l{\cel{3}DEFIRS.}
\l{}
\l{\cel{1}-\cel{12}\textquotedbl{}\cel{6}\sou{\textquotedbl{} \textquotedbl{}-}\cel{1}\textquotedbl{}\textquotedbl{}\cel{3}\textquotedbl{}\cel{2}\textquotedbl{} \textquotedbl{}-\textquotedbl{}- - \textquotedbl{}\cel{5}-- --\cel{2}--\textquotedbl{}\textquotedbl{}}
}
